\section{The Natural Number Systems}

The set $\set*{1, 2, 3, \dots}$ of \textbf{natural numbers} is denoted by $\N$, and is defined by the following axioms:
\begin{itemize}
    \item $1 \in \N$
    \item $\forall n \in \N, s(n) \in \N$
    \item $\nexists n \in \N, s(n) = 1$
    \item $\forall m, n \in \N, s(m) = s(n) \implies m = n$
    \item If $S \subseteq \N$ satisfies
    \begin{itemize}
        \item $1 \in S$
        \item $\forall n \in S, s(n) \in S$
    \end{itemize}
    then $S = \N$
\end{itemize}
where $s(n)$ is called the \textbf{successor} of $n$.
These axioms are called the \textbf{Peano axioms} (or \emph{Peano Postulates}).

\bigskip

\begin{remark}
    The \textbf{Peano axioms} do not guarantee the \emph{existence} of $\N$, which is guaranteed by the
    \textbf{axiom of infinity} in the \textbf{ZFC axioms}. In this textbook, the \textbf{ZFC axioms}
    is not discussed, so we will just assume that $\N$ exists.
\end{remark}

\bigskip

\begin{definition} \label{an:def:mathematical-induction}
    Let $P(n)$ be a statement about $n \in \N$. Define
    \[
    S \coloneqq \setbuilder*{n \in \N}{P(n) \text{ is true}} \subseteq \N
    \]
    If $P(n)$ satisfies
    \begin{itemize}
        \item $P(1)$ is true
        \item $\forall n \in \N, P(n)$ is true $\implies P(s(n))$ is true
    \end{itemize}
    Then $S$ satisfies
    \begin{itemize}
        \item $1 \in S$
        \item $\forall n \in S, s(n) \in S$
    \end{itemize}
    Therefore,
    \[
    \begin{array}{c}
        S = \N \\
        \Downarrow \\
        \forall n \in \N, P(n) \text{ is true}
    \end{array}
    \]
    This method of proof is called the \textbf{principle of mathematical induction}.
\end{definition}

\bigskip

\begin{definition} \label{an:def:operations-on-N}
    Let $n, m \in \N$ be given.
    \begin{itemize}
        \item The \textbf{addition} $+$ on $\N$ is defined as
        \begin{itemize}
            \item $n + 1 \coloneqq s(n)$
            \item $n + s(m) \coloneqq s(n + m)$
        \end{itemize}
        \item The \textbf{multiplication} $\cdot$ on $\N$ is defined as
        \begin{itemize}
            \item $n \cdot 1 \coloneqq n$
            \item $n \cdot s(m) \coloneqq n \cdot m + n$
        \end{itemize}
    \end{itemize}
\end{definition}

\bigskip

\begin{remark}
    Strictly speaking, \Cref{an:def:operations-on-N} requires the \textbf{recursion theorem} to be strictly defined.
    In this textbook, the \textbf{recursion theorem} is not discussed, so we will just assume that the operations
    on $\N$ are \textbf{well-defined}.
\end{remark}

\bigskip

\begin{theorem} \label{an:thm:properties-of-operations-on-N}
    Let $n, m, k \in \N$ be given. Then the following properties hold:
    \begin{itemize}
        \item $n + m = m + n$
        \item $(n + m) + k = n + (m + k)$
        \item $n \cdot m = m \cdot n$
        \item $(n \cdot m) \cdot k = n \cdot (m \cdot k)$
        \item $n \cdot (m + k) = n \cdot m + n \cdot k$
    \end{itemize}
\end{theorem}

\begin{proof}
    We will not discuss the proof of this theorem, since it relies on the recursive definitions.
\end{proof}

\bigskip

\begin{theorem} \label{an:thm:well-ordering-property}
    $\N$ has the \textbf{well-ordering property}, which states that
    \[
    \forall S \subseteq \N, S \neq \vnot \implies \exists m \in S, m = \min S
    \]
\end{theorem}

\begin{proof}
    Assume that
    \[
    \exists S \subseteq \N, S \neq \vnot \land (\nexists m \in S, m = \min S)
    \]
    Define
    \begin{itemize}
        \item $T \coloneqq \N \setminus S \subseteq \N$
        \item $L \coloneqq \setbuilder*{n \in \N}{\set*{1, 2, \dots, n} \subseteq T} \subseteq \N$
    \end{itemize}
    We will show that
    \[
    \begin{array}{c}
        L = \N \\
        \Downarrow \\
        T = \N \\
        \Downarrow \\
        S = \vnot \\
        \Downarrow \\
        \bot
    \end{array}
    \]
    \begin{itemize}
        \item \textbf{Base case:} If $1 \in S$, then
        \[
        \begin{array}{c}
            1 \in S \\
            \Downarrow \\
            1 = \min S \\
            \Downarrow \\
            \bot
        \end{array}
        \]
        Therefore,
        \[
        \begin{array}{c}
            1 \notin S \\
            \Downarrow \\
            1 \in T \\
            \Downarrow \\
            1 \in L
        \end{array}
        \]
        \item \textbf{Inductive step:} Let $n \in L$ be given. Then
        \[
        \begin{array}{c}
            n \in L \\
            \Downarrow \\
            \set*{1, 2, \dots, n} \subseteq T \\
            \Downarrow \\
            \set*{1, 2, \dots, n} \cap S = \vnot
        \end{array}
        \]
        If $n + 1 \in S$, then
        \[
        \begin{array}{c}
            n + 1 \in S \\
            \Downarrow \\
            n + 1 = \min S \\
            \Downarrow \\
            \bot
        \end{array}
        \]
        Therefore,
        \[
        \begin{array}{c}
            n + 1 \notin S \\
            \Downarrow \\
            n + 1 \in T \\
            \Downarrow \\
            \set*{1, 2, \dots, n + 1} \subseteq T \\
            \Downarrow \\
            n + 1 \in L
        \end{array}
        \]
    \end{itemize}
    By the \emph{principle of mathematical induction}, $L = \N$.
\end{proof}


